\documentclass[12pt]{amsart}
\pdfoutput=1
\usepackage{latexsym, amsmath, amssymb, amsthm, bm}
\usepackage{calrsfs}
\usepackage{subfig}

\usepackage[mathscr]{eucal}
\usepackage{mathtools}
\usepackage{paralist}
\usepackage[alphabetic]{amsrefs}
\usepackage[T1]{fontenc}
\usepackage{mathptmx}
\usepackage{microtype}
\usepackage{tikz}
\renewcommand{\baselinestretch}{1.05}
\usepackage[colorlinks=true,linkcolor=blue,urlcolor=blue, citecolor=blue]%
    {hyperref}
\linespread{1.06}
%%--LAYOUT--------------------------------------------------------------------
\usepackage[paperwidth=8.5in, paperheight=11in, inner=2.5cm, outer=2.5cm,% 
    top=2.5cm, bottom=2.5cm]{geometry}
%%--OTHER ENVIRONMENTS--------------------------------------------------------
\newtheorem{lemma}{Lemma}[section]
\newtheorem{theorem}[lemma]{Teorema}
\newtheorem{corollary}[lemma]{Corolario}
\newtheorem{proposition}[lemma]{Proposición}
\theoremstyle{definition}
\newtheorem{definition}[lemma]{Definición}
\newtheorem{remark}[lemma]{Observación}
\newtheorem{examplex}[lemma]{Ejemplo}
\newtheorem{procedure}[lemma]{Procedure}

\newenvironment{exm}
  {\pushQED{\qed}\renewcommand{\qedsymbol}{$\diamond$}\examplex}
  {\popQED\endexamplex}
\newtheorem{question}[lemma]{Question}
\newtheorem{conjecture}[lemma]{Conjecture}
\newtheorem{Construction}[lemma]{Construction}
\newtheorem{algorithm}[lemma]{Algorithm}
%\theoremstyle{remark}
\newcommand{\define}[1]{{\bfseries\itshape #1}}
%%--MATH----------------------------------------------------------------------
\newcommand{\relphantom}[1]{\mathrel{\phantom{#1}}}
\newcommand{\abs}[1]{\ensuremath{\left| #1 \right|}}
%\newcommand{\norm}[1]{\ensuremath{\left\| #1 \right\|}}
\newcommand{\ideal}[1]{\ensuremath{\left\langle #1 \right\rangle}}
\renewcommand{\geq}{\geqslant}
\renewcommand{\leq}{\leqslant}
\newcommand{\coloneq}{\ensuremath{\mathrel{\mathop :}=}}
\newcommand{\transpose}{\ensuremath{\textsf{\upshape T}}} 
%% Blackboard bolds symbols
\renewcommand{\AA}{\ensuremath{\mathbb{A}}} 
\newcommand{\CC}{\ensuremath{\mathbb{C}}} 
\newcommand{\kk}{\ensuremath{\Bbbk}} 
\newcommand{\NN}{\ensuremath{\mathbb{N}}}
\newcommand{\PP}{\ensuremath{\mathbb{P}}} 
\newcommand{\QQ}{\ensuremath{\mathbb{Q}}} 
\newcommand{\RR}{\ensuremath{\mathbb{R}}} 
\newcommand{\EE}{\ensuremath{\mathbb{E}}} 

\renewcommand{\P}{\ensuremath{P}} 
\newcommand{\Q}{\ensuremath{Q}} 

%\renewcommand{\SS}{\ensuremath{\mathbb{S}}} 
\newcommand{\ZZ}{\ensuremath{\mathbb{Z}}} 
%% Script symbols
\newcommand{\sC}{\ensuremath{\mathcal{C}}}
\newcommand{\sE}{\ensuremath{\mathcal{E}}}
\newcommand{\sF}{\ensuremath{\mathcal{F}}}
\newcommand{\calF}{\ensuremath{\mathcal{F}}}
\newcommand{\calL}{\ensuremath{\mathcal{L}}}

\newcommand{\sH}{\ensuremath{\mathcal{H}}}
\newcommand{\sL}{\ensuremath{\mathcal{L}}}
\newcommand{\sO}{\ensuremath{\mathcal{O}}}

\newcommand{\diam}{\operatorname{diam}}
\newcommand{\Osh}{{\mathcal O}}
\newcommand{\Ish}{{\mathcal I}}
\newcommand{\calI}{{\mathcal I}}

%% Operators
\DeclareMathOperator{\codim}{codim}
\DeclareMathOperator{\coker}{Coker}
\DeclareMathOperator{\conv}{conv}
\DeclareMathOperator{\Cox}{Cox}
\DeclareMathOperator{\depth}{depth}
\DeclareMathOperator{\gon}{gon}
\DeclareMathOperator{\Grass}{Gr}
\DeclareMathOperator{\green}{\textit{a}}
\DeclareMathOperator{\HH}{H}
\DeclareMathOperator{\HHH}{\mathbb{H}}
\DeclareMathOperator{\Hom}{Hom}
\DeclareMathOperator{\id}{id}
\DeclareMathOperator{\Image}{Im}
\DeclareMathOperator{\kept}{\lambda}
\DeclareMathOperator{\Ker}{Ker}
\DeclareMathOperator{\len}{\ell}
\DeclareMathOperator{\py}{py}
\DeclareMathOperator{\Pic}{Pic}
\DeclareMathOperator{\Proj}{Proj}
\DeclareMathOperator{\pr}{p}
\DeclareMathOperator{\qp}{qp}
\DeclareMathOperator{\qr}{qr}
\DeclareMathOperator{\rank}{rank}
\DeclareMathOperator{\reg}{reg}
\DeclareMathOperator{\Sos}{\Sigma}
\DeclareMathOperator{\Span}{Span}
\DeclareMathOperator{\Spec}{Spec}
\DeclareMathOperator{\sw}{sw}
\DeclareMathOperator{\Sym}{Sym}
\DeclareMathOperator{\topint}{int}
\DeclareMathOperator{\Tor}{Tor}
\DeclareMathOperator{\tw}{tw}
\DeclareMathOperator{\variety}{V}
%% Personalized comments
\newcommand\mv[1]{{\color{blue}(MV: #1)}}
\newcommand\mve[1]{{\color{red}(MV: #1)}}
\newcommand\jh[1]{{\color{red}(JH: #1)}}

\title{Guía de estudio II (Probabilidad y Estadística)}

\begin{document}

\maketitle
El objetivo de este documento es servir de guía para los estudiantes en su preparación para los parciales del curso.
Los enunciados marcados con Lema o Teorema son fundamentales y hay que conocer sus demostraciones. 

\section{Función generadora de momentos}

\begin{definition} Si $X$ es una variable aleatoria con valores reales definimos su función generadora de momentos mediante la fórmula
    \[m_X(t)=\EE[e^{tX}]\]
Decimos que la variable $X$ tiene una función generadora de momentos si existe un número $\epsilon>0$ tal que la función $m_X(t)$ existe (i.e. la integral converge) para todo $t\in (-\epsilon,\epsilon)$.
\end{definition}

La función generadora de momentos es una de las maneras en las que podemos representar una variable aleatoria. 
El siguiente Lema reúne algunas propiedades básicas de la función generadora de momentos (no es necesario conocer su demostración pues la demostración que vimos depende del Teorema de convergencia dominada y las ecuaciones de Cauchy-Riemann, ambos temas importantes pero por fuera de este curso).

\begin{lemma} Si la función generadora de momentos $m_X(t)$ existe para $t\in (-\epsilon,\epsilon)$ entonces,
    \begin{enumerate}
        \item La función $m_X(t)$ es continua en $(-\epsilon,\epsilon)$.
        \item La función $m_X(t)$ es diferenciable en $(-\epsilon,\epsilon)$.
        \item La función $m_X(t)$ es analítica en $(-\epsilon,\epsilon)$ y en particular:
        \begin{enumerate}
            \item Para todo entero positivo $k$, la $k$ésima derivada en cero $m_X^{(k)}(0) = \EE[X^k]$ 
            \item Para todo $t\in (-\epsilon,\epsilon)$ se tiene que
            \[m_X(t)=\sum_{k=0}^\infty \frac{\EE[X^k]}{k!}\,t^k \]
        \end{enumerate}
    \end{enumerate}
\end{lemma}

En particular, el Lema muestra cómo usar la función generadora de momentos para calcular los momentos de una variable aleatoria.

Adicionalmente el siguiente Lema (del que no es necesario conocer la demostración pero sí el enuncciado) aclara en qué sentido la función generadora de momentos $m_X(t)$ captura la distribución de $X$.

\begin{lemma} Las siguientes afirmaciones son verdaderas:
    \begin{enumerate}
        \item Si $X,Y$ son variables aleatorias cuyas funciones generadoras de momento coinciden en algún intervalo alrededor de $t=0$ (i.e. si $m_X(t)=m_Y(t)$ para todo $t\in (-\epsilon,\epsilon)$ para algún $\epsilon>0$) entonces las funciones de distribución acumulada de $X$ y de $Y$ coinciden.
        \item Sean $(X_n)_{n\in \mathbb{N}}$ y $Z$ variables aleatorias cuyas funciones generadoras de momento existen en algún intervalo $(-\epsilon,\epsilon)$. Si $\lim_{n\rightarrow \infty} m_{X_n}(t) = m_Z(t)$ para todo $t$ en algún intervalo alrededor del orgien entonces para todo $\beta\in\RR$ en la que la función de distribución acumulada de $Z$ sea continua se tiene que 
        \[\lim_{n\rightarrow \infty}\PP\left\{X_n\leq \beta\right\} = \PP\{Z\leq \beta\}\] 
    \end{enumerate}
    
\end{lemma}


\begin{examplex} \label{ex: Gauss} Si $X$ es una variable aleatoria con distribución gaussiana, media $\mu$ y varianza $\sigma^2$, entonces
\[
m_X(t)=e^{\mu t+\frac{1}{2}\sigma^2 t^2}\qquad\text{para todo }t\in\RR.
\]
En efecto, como la densidad de $X$ es
\[
f(x)=\frac{1}{\sigma\sqrt{2\pi}}\,\exp\!\left(-\frac{(x-\mu)^2}{2\sigma^2}\right),
\]
se tiene
\[
m_X(t)=\int_{-\infty}^{\infty}e^{tx}f(x)\,dx
   =\frac{1}{\sigma\sqrt{2\pi}}\int_{-\infty}^{\infty}\exp\!\left(tx-\frac{(x-\mu)^2}{2\sigma^2}\right)\,dx.
\]
Completando el cuadrado en el exponente,
\[
tx-\frac{(x-\mu)^2}{2\sigma^2}
   =-\frac{\left(x-(\mu+\sigma^2t)\right)^2}{2\sigma^2}+\mu t+\frac{1}{2}\sigma^2t^2,
\]
de modo que
\[
m_X(t)=e^{\mu t+\frac{1}{2}\sigma^2t^2}\int_{-\infty}^{\infty}\frac{1}{\sigma\sqrt{2\pi}}
   \exp\!\left(-\frac{\left(x-(\mu+\sigma^2t)\right)^2}{2\sigma^2}\right)\,dx.
\]
Con el cambio de variables $y=\dfrac{x-(\mu+\sigma^2t)}{\sigma}$ (de modo que $dx=\sigma\,dy$), la integral se convierte en
\[
\int_{-\infty}^{\infty}\frac{1}{\sqrt{2\pi}}e^{-y^2/2}\,dy=1,
\]
En consecuencia $m_X(t)=e^{\mu t+\frac{1}{2}\sigma^2t^2}\qquad\text{para todo }t\in\RR$.
De la expansión de Taylor de $m_X(t)$ alrededor de $t=0$ obtenemos simultáneamente todos los momentos de $X$. Para simplificar, asumamos que $\mu=0$; entonces $m_X(t)=e^{\frac{1}{2}\sigma^2 t^2}$ y su expansión de Taylor en $t=0$ es
\[
m_X(t)=\sum_{k=0}^{\infty}\frac{\left(\frac{1}{2}\sigma^2 t^2\right)^k}{k!}
   =\sum_{k=0}^{\infty}\frac{(\sigma^2/2)^k}{k!}\,t^{2k}.
\]
Como en esta expansión solo aparecen potencias pares de $t$, el coeficiente de $t^n$ es
\[
\begin{cases}
\dfrac{(\sigma^2/2)^{n/2}}{(n/2)!} & \text{si }n\text{ es par},\\[6pt]
0 & \text{si }n\text{ es impar}.
\end{cases}
\]
Pero por el Lema anterior el coeficiente de $t^n$ en la expansión de Taylor de $m_X(t)$ es $\dfrac{\EE[X^n]}{n!}$; igualando ambos coeficientes obtenemos la fórmula para los momentos de una gaussiana centrada:
\[
\EE[X^n]=\begin{cases}
\dfrac{n!}{(n/2)!}\left(\dfrac{\sigma^2}{2}\right)^{n/2}=(n-1)!!\,\sigma^n & \text{si }n\text{ es par},\\[8pt]
0 & \text{si }n\text{ es impar},
\end{cases}
\]
donde $(2k-1)!!=1\cdot 3\cdot 5\cdots(2k-1)$ denota el doble factorial. En particular,
\[
\EE[X]=0,\qquad \EE[X^2]=\sigma^2,\qquad \EE[X^3]=0,\qquad \EE[X^4]=3\sigma^4,\qquad \EE[X^6]=15\sigma^6.
\]
\end{examplex}    

\begin{examplex} 
    Si $X$ es una variable aleatoria con distribución uniforme en el intervalo $[a,b]$, $a<b$, entonces
\[
m_X(t)=\begin{cases}
\dfrac{e^{tb}-e^{ta}}{t(b-a)} & \text{si }t\neq 0,\\[6pt]
1 & \text{si }t=0.
\end{cases}
\]
En efecto, la densidad de $X$ es $f(x)=\dfrac{1}{b-a}$ para $a\le x\le b$ (y $0$ en otro caso), por lo que
\[
m_X(t)=\int_{-\infty}^{\infty}e^{tx}f(x)\,dx
   =\frac{1}{b-a}\int_{a}^{b}e^{tx}\,dx
   =\frac{1}{b-a}\cdot\frac{e^{tb}-e^{ta}}{t}
   =\frac{e^{tb}-e^{ta}}{t(b-a)}
\]
para $t\neq 0$ (y $m_X(0)=1$, que es el límite cuando $t\to 0$).

De la expansión de Taylor de $m_X(t)$ alrededor de $t=0$ obtenemos todos los momentos de $X$. Al expandir el numerador,
\[
e^{tb}-e^{ta}=\sum_{k=0}^{\infty}\frac{t^k}{k!}\left(b^k-a^k\right)
   =t\sum_{n=0}^{\infty}\frac{t^n}{(n+1)!}\left(b^{n+1}-a^{n+1}\right),
\]
de modo que
\[
m_X(t)=\frac{1}{b-a}\sum_{n=0}^{\infty}\frac{b^{n+1}-a^{n+1}}{(n+1)!}\,t^n.
\]
Como por el Lema anterior el coeficiente de $t^n$ en la expansión de Taylor es $\dfrac{\EE[X^n]}{n!}$, igualando ambos coeficientes obtenemos la fórmula para los momentos de la uniforme:
\[
\EE[X^n]=\frac{n!}{(n+1)!}\cdot\frac{b^{n+1}-a^{n+1}}{b-a}
   =\frac{b^{n+1}-a^{n+1}}{(n+1)(b-a)}.
\]
En particular,
\[
\EE[X]=\frac{a+b}{2},\qquad
\EE[X^2]=\frac{a^2+ab+b^2}{3},\qquad
\EE[X^3]=\frac{a^3+a^2b+ab^2+b^3}{4}.
\]
\end{examplex}

\subsection{Aplicaciones de la función generadora de momentos}

La primera aplicación es uno de los teoremas más útiles de la probabilidad. Explica la ubicuidad de la distribución normal pues captura cómo son las fluctuaciones de la media muestral $\overline{X_n}$ alrededor de la media. La distribución Gaussiana es una descripción universal de las fluctuaciones que resultan de agregar muchas contribuciones aleatorias \emph{de fuentes independientes}.

\begin{theorem} {\it Teorema del Límite Central} Suponga que $X_1,\dots, X_n$ son variables aleatorias
independientes e idénticamente distribuidas, con media $\mu$ y varianza $\sigma^2$, y que
$m_{X_j}(t)$ existe en algún intervalo alrededor de cero. Si $\overline{X_n}:=\frac{1}{n}\sum_{j=1}^n X_j$,
entonces 
\[\frac{\overline{X_n}-\EE[\overline{X_n}]}{\sqrt{{\rm Var}(\overline{X_n})}} = \frac{\overline{X_n}-\mu}{\frac{\sigma}{\sqrt{n}}}\]
y para todo $\beta\in \RR$
\[\lim_{n\rightarrow \infty}\PP\left\{ \frac{\overline{X_n}-\mu}{\sigma/\sqrt{n}}\leq \beta\right\}=\Phi(\beta),
\]
donde $\Phi(z)$ es la función de distribución acumulada de la normal estándar $\mathcal{N}(0,1)$.
\end{theorem}
\begin{proof} Sin pérdida de generalidad asumamos que las $X_i$ tienen media $0$ (restando $\mu$ a cada una y redefiniéndolas, si es necesario). Como las $X_i$ son independientes, el valor esperado de un producto es el producto de los
valores esperados, y
\[\EE\left[e^{\frac{t\sqrt{n}}{\sigma}\,\overline{X_n}}\right]
 =\EE\left[e^{\frac{t}{\sigma\sqrt{n}}\sum_{j=1}^n X_j}\right]
 =\prod_{j=1}^n m_{X_j}\left(\frac{t}{\sigma\sqrt{n}}\right)
 =\left(m_{X_1}\left(\frac{t}{\sigma\sqrt{n}}\right)\right)^{n}.
\]
Como $m_{X_1}$ es analítica en cero y $\EE[X_1]=0$ (no hay término lineal), su expansión de Taylor en el
origen da
\[
m_{X_1}\left(\frac{t}{\sigma\sqrt{n}}\right)=1+\frac{t^2}{2n}+O\!\left(\frac{1}{n^{3/2}}\right),
\]
de donde
\[
\left(m_{X_1}\left(\frac{t}{\sigma\sqrt{n}}\right)\right)^n
=\exp\!\left(n\log\left(1+\frac{t^2}{2n}+O\!\left(\frac{1}{n^{3/2}}\right)\right)\right).
\]
Puesto que $n\log\left(1+\frac{t^2}{2n}+O\!\left(\frac{1}{n^{3/2}}\right)\right)\to \frac{t^2}{2}$ cuando
$n\to\infty$, se sigue que
\[
\lim_{n\to\infty}\left(m_{X_1}\left(\frac{t}{\sigma\sqrt{n}}\right)\right)^n = e^{t^2/2},
\]
que, por el ejemplo~\ref{ex: Gauss}, es la función generadora de momentos de la normal estándar.
\end{proof}

\section{Estadística paramétrica}

El objetivo de la estadística es entender el comportamiento de una variable de interés (típicamente numérica o categórica) en una población.
La hipótesis fundamental de la \emph{estadística paramétrica} es que esta cantidad es aleatoria y tiene una \emph{distribución poblacional} fija. 
Asumimos:
\begin{enumerate}
    \item Que conocemos una familia de distribuciones de probabilidad (llamada \emph{modelo paramétrico}) que depende de un \emph{vector de parámetros $\theta$} cuando $\theta$ recorre un conjunto de parámetros permitidos $\Theta$. 
\item Que la distribución poblacional se obtiene cuando los parámetros toman un valor específico  $\theta=\theta^*$ llamado el \emph{valor verdadero}
\end{enumerate}

Nuestro objetivo es intentar estimar el valor verdadero de $\theta^*$ a partir del análisis de la cantidad en una muestra de tamaño $n$ extraída de la población. Por el carácter aleatorio de la muestra y de la cantidad de interés, nuestra estimación será necesariamente aleatoria. Nuestro propósito será entonces garantizar que, con alta probabilidad, nuestras estimaciones estén cerca (en un sentido a definir) del valor verdadero $\theta^*$ de los parámetros. Algunos ejemplos aclararán el punto de vista paramétrico,

\begin{examplex}
    Queremos entender la altura $Z$ de las personas en una población. Este problema permite muchos modelos paramétricos posibles. La selección del modelo paramétrico es nuestra \emph{hipótesis fundamental} y la que requiere más escrutinio (postular un modelo es más arte que ciencia y típicamente parte de observar los datos disponibles con suficiente atención). Para ilustrar las posibilidades podríamos considerar los siguientes modelos sencillos:
    \begin{enumerate}
        \item La altura tiene distribución normal con media desconocida $\mu$ y varianza conocida $\sigma^2=10^2$ en centímetros (recuerden que, para la distribución normal con varianza $\sigma^2$ el intervalo de $\mu \pm 3\sigma$ debe contener más del $99\%$ de los datos, lo cual en el caso de la altura se traduce a un rango de $\pm 30$ cm con la selección de varianza de arriba). En este caso $\theta=\mu$ y la familia paramétrica sería 
        \[f_{\theta}(x)=\frac{1}{\sqrt{2\pi}} e^{-\frac{(x-\theta)^2}{2\cdot 10^2}}\]
        La \emph{distribución poblacional} sería la normal de arriba para algún valor $\theta^*=\mu$ fijo pero desconocido para nosotros.

        \item La altura tiene distribución normal con media $\mu$ y varianza $\sigma^2$, ambas desconocidas, es decir $\theta=(\mu,\sigma^2)$.
        La familia paramétrica sería:
        \[f_{\theta}(x)=\frac{1}{\sqrt{2\pi\sigma^2}} e^{-\frac{(x-\mu)^2}{2\sigma^2}},\qquad \theta=(\mu,\sigma^2)\in\RR\times(0,\infty)\]
        La \emph{distribución poblacional} sería la normal de arriba para algún valor $\theta^*=(\mu,\sigma^2)$ fijo pero desconocido para nosotros.
        
        \item La altura tiene una distribución descrita con una mezcla de dos distribuciones normales equiprobables con parámetros $(\mu_1,\sigma_1^2)$ y $(\mu_2,\sigma_2^2)$ que representan las alturas de mujeres y hombres en la población. En ese caso $\theta=\left(\mu_1,\sigma_1^2,\mu_2,\sigma_2^2\right)$.       
        La familia paramétrica sería:
        \[f_{\theta}(x)=\frac{1}{2}\cdot\frac{1}{\sqrt{2\pi\sigma_1^2}} e^{-\frac{(x-\mu_1)^2}{2\sigma_1^2}}+\frac{1}{2}\cdot\frac{1}{\sqrt{2\pi\sigma_2^2}} e^{-\frac{(x-\mu_2)^2}{2\sigma_2^2}},\qquad \theta\in\RR\times(0,\infty)\times\RR\times(0,\infty)\]
        La \emph{distribución poblacional} sería esta mezcla de dos normales para algún valor $\theta^*$ fijo pero desconocido para nosotros, donde $\frac{1}{2}$ es el peso de cada componente al ser equiprobables.
    \end{enumerate}
\end{examplex}

Disponemos de una muestra de tamaño $n$ de la cantidad de interés (las alturas de $n=100$ personas en la población, por ejemplo). Asumimos que la muestra es una realización de variables aleatorias $X_1,\dots, X_n$ que son \emph{independientes} y que tienen la distribución poblacional (es decir la distribución paramétrica cuando $\theta=\theta^*$).

\begin{definition} Un estimador insesgado de una función  $\tau(\theta)$ de los parámetros de la familia es una función $T(X_1,\dots, X_n)$ de la muestra que cumple 
    \[\forall \theta \left(\EE_\theta[T(X_1,\dots, X_n)]= \tau(\theta)\right) \]
    donde $\EE_\theta[\bullet]$ significa que el valor esperado esta siendo calculado asumiendo que las $X_i$ tienen la distribución resultante de asumir que la distribución poblacional esta la dada por el parámetro $\theta$.    
\end{definition}

El siguiente Teorema es muy importante pues provee algunos ejemplos básicos de estimadores insesgados. La calidad de un estimador insesgado depende de su distribución y en particular de qué tan concentrada se encuentre alrededor de su media, así que es útil conocer las distribuciones de estimadores como se hace en la parte $(3)$.

\begin{theorem} Sean $X_1,\dots, X_n$ variables aleatorias independientes con media $\mu$ y varianza $\sigma^2$. Entonces las siguientes propiedades se cumplen:
\begin{enumerate}
    \item $\overline{X_n}$ es un estimador insesgado para $\mu$.
    \item $S^2:=\frac{1}{n-1}\sum_{j=1}^n \left(X_j-\overline{X_n}\right)^{2}$ es un estimador insesgado para $\sigma^2$. 
    \item Si todas las $X_j$ son variables aleatorias con distribución normal $\mathcal{N}(\mu,\sigma^2)$ entonces:
    \begin{enumerate}
        \item $\overline{X_n}$ tiene distribución normal $\mathcal{N}(\mu,\frac{\sigma^2}{n})$
        \item $\frac{(n-1)S^2}{\sigma^2}$ tiene distribución $\chi^2$ con $n-1$ grados de libertad
        \item El estadístico $W:=\frac{\overline{X_n}-\mu}{S/\sqrt{n}}$ tiene distribución t de Student con $n-1$ grados de libertad.
    \end{enumerate} 
\end{enumerate}
\end{theorem}

\subsection{Intervalos de confianza y pruebas de hipótesis}

Nuestro estimador $T(X_1,\dots, X_n)$ del parámetro real $\theta^*$ es una función de las muestras y, por lo tanto, es una variable aleatoria. Cuando evaluamos $T$ en una realización de la muestra obtenemos un número y nos gustaría cuantificar qué tan lejos esta ese número del valor real $\theta_0$, o al menos garantizar que esta cerca con una probabilidad razonable. Los intervalos de confianza son una manera concreta de realizar esta idea. Más precisamente,

\begin{definition} Dado un número $\alpha\in (0,1)$, decimos que dos estimadores $L(X_1,\dots, X_n)$ y $U(X_1,\dots, X_n)$ definen un \emph{intervalo de confianza} para el parámetro $\beta$ si sucede que
    \[\PP\left\{ L(X_1,\dots, X_n)\leq \beta \leq U(X_1,\dots, X_n)\right\}=1-\alpha\]
donde la probabilidad se calcula asumiendo que la muestra $X_i$ es una sucesión independiente que tiene la distribución de nuestro modelo para algún valor $\theta=\theta_0$ de los parámetros.    
\end{definition}






\end{document}

